A dolgozat elkészítése során generatív mesterséges intelligencia eszközt
(Anthropic Claude) \textbf{segédeszközként} vettem igénybe. A főbb felhasználási
területek a következők voltak:

\begin{itemize}
  \item a forráskód hibakeresése és egyes kódrészletek működésének tisztázása;
  \item az ábrák egy részét előállító \texttt{matplotlib}-szkriptek elkészítése
  (a sematikus diagramok és az eredmény-grafikonok);
  \item a \LaTeX-dokumentum formázása és szerkesztése;
  \item a szöveg nyelvi és stilisztikai csiszolása (helyesírás, terminológiai
  következetesség);
  \item a szakirodalmi hivatkozások formai és tartalmi ellenőrzése;
  \item a kivonat angol és román nyelvű fordításának előkészítése.
\end{itemize}

Az eszköz által javasolt tartalmat a végleges változatba illesztés előtt minden
esetben ellenőriztem.
